\PassOptionsToPackage{quiet}{xeCJK}
\documentclass[withoutpreface,bwprint]{format}
\usepackage{etoolbox}
\usepackage{ctex}
\BeforeBeginEnvironment{tabular}{\zihao{-5}}
% 参考文献改用纯tex格式，不再需要natbib+bibtex
\usepackage[framemethod=TikZ]{mdframed}
\usepackage{url}
\usepackage{array}
\usepackage{tabularx}
\usepackage{longtable}
\newcolumntype{C}{>{\centering\arraybackslash}X}
\newcolumntype{R}{>{\raggedleft\arraybackslash}X}
\newcolumntype{L}{>{\raggedright\arraybackslash}X}

\renewcommand{\textbf}[1]{{\song\bfseries #1}}

\usepackage{tikz}
\usetikzlibrary{arrows.meta}
\tikzset{
  box/.style={rectangle, draw, minimum width=3.2cm, minimum height=0.9cm, align=center},
  arrow/.style={thick, -{Stealth}}
}

% 数值宏由 results_ledger.json 通过 ledger.py --emit-tex 生成。
% 缺失时保留可见失败标记，不能静默编译成看似完整的论文。
\IfFileExists{numbers.tex}{%
  \input{numbers.tex}%
}{%
  \newcommand{\NUMBERSMISSING}{\textbf{[NUMBERS-MISSING]}}%
  \typeout{[NUMBERS-MISSING] numbers.tex 未找到：先运行 ledger.py --emit-tex}%
}








\title{论文标题}



%%%%%%%%%%%%%%%%%%%%%%%%%%%%%%%%%%%%%%%%%%%%%%%%%%%%%%%%%%%%%
%% 正文
\begin{document}

% 标题
{\centering\zihao{3}\bfseries 论文标题\par}

\setcounter{page}{1}

\thispagestyle{empty}

% ── 全文通用形态（Mrite 高教社杯规范）──
% 正文一律自然段；禁分点符号、禁 \textbf（摘要“针对问题X”与问题重述“问题N：”除外）。
% 跨页表用 longtable，短表用 tabularx；列宽比例总和 = 1.04 − 0.04 × N（N 为列数），格内文字不换行。
% 图宽 0.8\textwidth，表宽 \textwidth；图题交给 \caption，绘图代码不写 set_title()；图表必须在正文被引用。
% 章与章、节与节之间一律不加 \newpage。
% ── 本节合同 ──
% 总字数 ≤900，开头段 ≤30 字，结尾段 ≤20 字；编译后必须恰为 1 页（\label{abstract:end} 供 .aux 程序化核验，不得删）。
% 每问一段，段内五动作固定：段首定位 → 建模主张 → 方法链条 → 结果句 → 验证句。
% 加粗只允许三类、每类至多一处：建立**XX模型** / 用**XX算法** / 得出**XXX数字或结论**。其余一律不加粗。
% 详细 move 库与 12 篇国奖语料出处见 references/abstract-moves.md（写摘要前读）。
%
\begin{abstract}

% ==== 开头段（2句话）====
...是...领域的重要研究课题。本文基于...数据，围绕...展开系统研究。

% ==== 逐问摘要 ====
\textbf{对于问题一,} 首先对...进行了...分析，接着采用...方法构建了...模型，最后通过...得到...。结果表明，...（关键数值）。

\textbf{对于问题二,} 首先...，接着...，最后...。结果表明...。

\textbf{对于问题三,} 首先...，接着...，最后...。结果表明...。

\textbf{对于问题四,} 首先...，接着...，最后...。结果表明...。

% ==== 结尾段（2句话）====
本文完成了...的研究，各模型性能优越、结论可靠。研究成果可为...提供理论依据和决策参考。


\keywords{关键词1；关键词2；关键词3；关键词4；关键词5}

\end{abstract}

\thispagestyle{empty}

\tableofcontents
\thispagestyle{empty}
\newpage

\setcounter{page}{1}

% ── 全文通用形态（Mrite 高教社杯规范）──
% 正文一律自然段；禁分点符号、禁 \textbf（摘要“针对问题X”与问题重述“问题N：”除外）。
% 跨页表用 longtable，短表用 tabularx；列宽比例总和 = 1.04 − 0.04 × N（N 为列数），格内文字不换行。
% 图宽 0.8\textwidth，表宽 \textwidth；图题交给 \caption，绘图代码不写 set_title()；图表必须在正文被引用。
% 章与章、节与节之间一律不加 \newpage。
% ── 本节合同 ──
% 问题背景 2–3 段：宏观背景 → 问题来源 → 当前困境 → 解决渴求，不分点。
% 问题重述用 \textbf{问题N：} 起头，每问 ≤2 行，只写问题本质，删背景。
%
\section{引言}
\subsection{问题背景}

% 第1段：宏观背景 → 第2段：问题来源与困境 → 第3段：解决渴求


\subsection{问题重述}

\textbf{问题1：}...

\textbf{问题2：}...

\textbf{问题3：}...

\textbf{问题4：}...


% ── 全文通用形态（Mrite 高教社杯规范）──
% 正文一律自然段；禁分点符号、禁 	extbf（摘要“针对问题X”与问题重述“问题N：”除外）。
% 跨页表用 longtable，短表用 tabularx；列宽比例总和 = 1.04 − 0.04 × N（N 为列数），格内文字不换行。
% 图宽 0.8	extwidth，表宽 	extwidth；图题交给 \caption，绘图代码不写 set_title()；图表必须在正文被引用。
% 章与章、节与节之间一律不加 
ewpage。
% ── 本节合同 ──
% 三段式：① 一句话总体表达；② 逐问串联“针对问题X建立了XX来解决XX”；③ 各问之间的递进关系。
% 本节不画流程图（流程图属于各问的“分析与准备”）。
%
\section{总体分析}

% 总述问题递进逻辑（纯文字，不画流程图）


% ── 全文通用形态（Mrite 高教社杯规范）──
% 正文一律自然段；禁分点符号、禁 \textbf（摘要“针对问题X”与问题重述“问题N：”除外）。
% 本节合同显式要求的结构（如 itemize、\textbf 或分点）属于模板既定例外，以本节合同为准。
% 跨页表用 longtable，短表用 tabularx；列宽比例总和 = 1.04 − 0.04 × N（N 为列数），格内文字不换行。
% 图宽 0.8\textwidth，表宽 \textwidth；图题交给 \caption，绘图代码不写 set_title()；图表必须在正文被引用。
% 章与章、节与节之间一律不加 \newpage。
% ── 本节合同 ──
% 两段式：第 1 段“为简化问题，本文做出以下假设：”，随后 itemize，每条 \textbf{假设N：} 且 ≤2 行，条数 = 问题数。
% 每条假设必须在模型检验章有对应处理（夹逼／对照／声明理由），否则不得写入。
%
\section{模型假设}

为简化问题，本文做出以下假设：

\begin{itemize}[itemindent=2em]
\item \textbf{假设1：}...
\item \textbf{假设2：}...
\item \textbf{假设3：}...
\item \textbf{假设4：}...
\end{itemize}


% ── 全文通用形态（Mrite 高教社杯规范）──
% 正文一律自然段；禁分点符号、禁 	extbf（摘要“针对问题X”与问题重述“问题N：”除外）。
% 跨页表用 longtable，短表用 tabularx；列宽比例总和 = 1.04 − 0.04 × N（N 为列数），格内文字不换行。
% 图宽 0.8	extwidth，表宽 	extwidth；图题交给 \caption，绘图代码不写 set_title()；图表必须在正文被引用。
% 章与章、节与节之间一律不加 
ewpage。
% ── 本节合同 ──
% 本节的 longtable 写法是全文表格的样板，其余各章一律沿用（含 \endfirsthead/\endhead/\endfoot 固定结构）。
% 2 列取 0.18+0.78；列宽比例总和必须等于 1.04 − 0.04 × N。
%
\section{符号说明}

本文使用的主要符号及其含义如表\ref{tab:符号说明}所示。

% longtable：可跨页，>{\centering\arraybackslash}p{} 居中，比例按内容调整
\begin{longtable}{>{\centering\arraybackslash}p{0.18\textwidth}>{\centering\arraybackslash}p{0.78\textwidth}}
  \caption{符号说明}
  \label{tab:符号说明} \\
  \toprule
  符号 & 说明 \\
  \midrule
  \endfirsthead
  \caption{符号说明（续）} \\
  \toprule
  符号 & 说明 \\
  \midrule
  \endhead
  \bottomrule
  \endfoot
  ... & ... \\
  ... & ... \\
  ... & ...
\end{longtable}


% ── 全文通用形态（Mrite 高教社杯规范）──
% 正文一律自然段；禁分点符号、禁 \textbf（摘要“针对问题X”与问题重述“问题N：”除外）。
% 跨页表用 longtable，短表用 tabularx；列宽比例总和 = 1.04 − 0.04 × N（N 为列数），格内文字不换行。
% 图宽 0.8\textwidth，表宽 \textwidth；图题交给 \caption，绘图代码不写 set_title()；图表必须在正文被引用。
% 章与章、节与节之间一律不加 \newpage。
% ── 本节合同 ──
% 只做装配：按问题数 N 增减 \input{5.X.问题X的建立求解.tex}，本文件不写正文。
%
\section{模型的建立与求解}

% ── 问题一模板 ──
% 本文件负责 5.1 标题、问题一流程图和 5.1.x 分节装配。
% 流程图不设置 \subsubsection，直接位于“问题一的模型的建立和求解”标题之后。
% 问题二至题面实际第 n 问复制本文件后，统一修改标题、标签、文字和 \input 路径。

\subsection{问题一的模型的建立和求解}

\begin{figure}[H]
  \centering
  \begin{tikzpicture}[node distance=0.7cm and 0.55cm]
    \node[flowbox] (n11) {数据理解};
    \node[flowbox, right=of n11] (n12) {数据预处理};
    \node[flowbox, right=of n12] (n13) {变量关系};
    \node[flowbox, below=of n11] (n21) {模型假设};
    \node[flowbox, below=of n12] (n22) {模型建立};
    \node[flowbox, below=of n13] (n23) {参数求解};
    \node[flowbox, below=of n21] (n31) {模型检验};
    \node[flowbox, below=of n22] (n32) {结果分析};
    \node[flowbox, below=of n23] (n33) {结论输出};
    \draw[arrow] (n11) -- (n12);
    \draw[arrow] (n12) -- (n13);
    \draw[arrow] (n13) -- (n23);
    \draw[arrow] (n23) -- (n22);
    \draw[arrow] (n22) -- (n21);
    \draw[arrow] (n21) -- (n31);
    \draw[arrow] (n31) -- (n32);
    \draw[arrow] (n32) -- (n33);
  \end{tikzpicture}
  \caption{问题一建模思路图}
  \label{fig:problem1_flow}
\end{figure}

% ── 问题一：5.1.1 模型准备 ──
% 问题分析已迁移至第2章；问题流程图直接位于 5.1 标题下，不在本文件重复。
% 本文件只承载模型准备。问题二至题面实际第 n 问复制本文件并改编号、标签和问题文字。
% 数据预处理（仅问题一）五段：原因 → 数据理解与提取 → 处理步骤（含公式）→ 前后统计量对比 → 总结。
% 算法选择五段：问题本质定性 → 候选算法与评估维度 → 对比表（写真实算法名与等级）→ 分析对比结论 → 总结引出。

\subsubsection{模型准备}

在进行建模之前，需要……。本节将从数据预处理和算法选择两方面进行准备。

原始数据来源于……，包含……条记录和……个字段。经检查发现以下问题：……（缺失值、量纲不统一、异常值或编码问题）。

针对上述问题，按以下步骤进行预处理：第一步，……（缺失值处理及参数）；第二步，……（标准化或归一化，并写出公式）；第三步，……（特征编码或衍生）。

预处理后，数据规模为……行×……列，各指标统计量如下表所示。

\begin{table}[H]
  \centering
  \caption{数据预处理前后统计量对比}
  \begin{tabularx}{\textwidth}{CCCCC}
    \toprule
    指标 & 预处理前均值 & 预处理前标准差 & 预处理后均值 & 预处理后标准差 \\
    \midrule
    …… & …… & …… & …… & …… \\
    \bottomrule
  \end{tabularx}
  \label{tab:preprocess_compare1}
\end{table}

综上所述，数据预处理完成，为后续建模提供了清洁的输入数据。

该问题本质上是一个……问题，需要在……条件下……。常用求解算法包括 A 算法、B 算法和 C 算法等。本节从……、……和……三方面进行评估。

\begin{table}[H]
  \centering
  \caption{算法能力对比}
  \begin{tabularx}{\textwidth}{>{\hsize=2.5\hsize\centering\arraybackslash}X>{\hsize=0.6\hsize\centering\arraybackslash}X>{\hsize=0.6\hsize\centering\arraybackslash}X>{\hsize=0.6\hsize\centering\arraybackslash}X>{\hsize=0.7\hsize\centering\arraybackslash}X}
    \toprule
    算法名称 & 维度1 & 维度2 & 维度3 & 稳定性 \\
    \midrule
    算法 A（缩写） & 优 & 良 & 中 & 良 \\
    算法 B（缩写） & 良 & 优 & 良 & 优 \\
    算法 C（缩写） & 中 & 中 & 优 & 中 \\
    \bottomrule
  \end{tabularx}
  \label{tab:algorithm_compare1}
\end{table}

通过对比可以看出，算法 A 在……方面表现最优，……。因此本文选择算法 A 作为核心求解方法。

综上所述，本节完成了数据预处理和算法选型，接下来将建立数学模型并进行求解。


\input{5.1.2.建模与求解.tex}


% 如有更多问题，按以下格式新增：
% \input{5.2.问题2的建立求解.tex}
% \input{5.3.问题3的建立求解.tex}


% ── 全文通用形态（Mrite 高教社杯规范）──
% 正文一律自然段；禁分点符号、禁 \textbf（摘要“针对问题X”与问题重述“问题N：”除外）。
% 跨页表用 longtable，短表用 tabularx；列宽比例总和 = 1.04 − 0.04 × N（N 为列数），格内文字不换行。
% 图宽 0.8\textwidth，表宽 \textwidth；图题交给 \caption，绘图代码不写 set_title()；图表必须在正文被引用。
% 章与章、节与节之间一律不加 \newpage。
% ── 本节合同（六小节缺一不可）──
% 本章须独立成章，不得把检验散落在各问结尾。对标同题人工基线，其检验章体量为全文第二大。
% ① 双路互证：两条独立实现算同一量，给一致性证据与差异定位。
% ② 与可核事实对表：题面或附件能验证的量，逐条列“声称值 ↔ 实测值”。
% ③ 参数灵敏度：逐参数扰动 → 结论是否翻转，给数字。
% ④ 样本量与收敛：证书随样本量的变化，说明为何足以分辨相邻档。
% ⑤ 稳健性与对照实验：换口径 / 换种子 / 换实现。
% ⑥ 适用边界与失效域。
%
\section{模型的检验}
\label{sec:validation}

为验证本文所建模型的有效性、可分辨性与适用边界，本章从六个方面独立检验。

\subsection{双路互证}

对每问的最终答案，用主路线与独立核验路线各算一次。两路须在实现层面独立（不同估计量、解析极限、暴力小实例或不同离散化），不得共享核心判定代码。差异超过分辨率 $\Delta$ 时必须定位原因，不取平均、不择优报喜。

\begin{table}[H]
  \centering
  \caption{双路互证结果}
  \label{tab:crosscheck}
  \begin{tabularx}{\textwidth}{CCCCC}
    \toprule
    问题 & 主路线取值 & 核验路线取值 & 差异 & 差异$/\Delta$ \\
    \midrule
    ... & ... & ... & ... & ... \\
    \bottomrule
  \end{tabularx}
\end{table}

如表\ref{tab:crosscheck}所示，...

\subsection{与可核事实对表}

题面与附件中一切可独立核验的量（计数、规模、边界值、给定档位），逐条列出声称值与实测值。此处对表的是\textbf{判定层与生成层两条链路}：生成器产出的构型统计量也须与附件同量级，只对其中一层等于裸奔。

\begin{table}[H]
  \centering
  \caption{可核事实对表}
  \label{tab:factcheck}
  \begin{tabularx}{\textwidth}{LCCC}
    \toprule
    题面/附件事实 & 出处 & 本文实测 & 是否一致 \\
    \midrule
    ... & ... & ... & ... \\
    \bottomrule
  \end{tabularx}
\end{table}

\subsection{参数灵敏度}

逐参数在工作点两侧扰动，报告输出变化量与结论是否翻转。灵敏度必须给数字，不得以“影响可忽略”定性结案。

\begin{table}[H]
  \centering
  \caption{参数灵敏度分析}
  \label{tab:sensitivity}
  \begin{tabularx}{\textwidth}{LCCCC}
    \toprule
    参数 & 扰动幅度 & 输出变化 & 变化$/\Delta$ & 结论是否翻转 \\
    \midrule
    ... & ... & ... & ... & ... \\
    \bottomrule
  \end{tabularx}
\end{table}

如表\ref{tab:sensitivity}所示，...

\subsection{样本量与收敛}

给出证书随样本量的变化曲线，并说明当前样本量为何足以分辨相邻档位。此处须区分两件事：\textbf{达标}（下界越过阈值）与\textbf{已分辨}（不确定度 $u \le \Delta/3$）。未分辨的结论须降到题面精度的下一档报告，不得报满精度。

% 收敛曲线图：有图后取消下面的注释，路径指向 求解/问题X/图片/
% \begin{figure}[H]
%   \centering
%   \includegraphics[width=0.8\textwidth]{../求解/图/收敛曲线.png}
%   \caption{证书下界随样本量的收敛}
%   \label{fig:convergence}
% \end{figure}

% 如图\ref{fig:convergence}所示，...

\subsection{稳健性与对照实验}

分别更换口径、更换随机种子、更换实现，检查结论是否稳定。答案稳定性复核须用至少两个独立主种子重跑终验，按题目精度取整后答案不变。

\begin{table}[H]
  \centering
  \caption{稳健性对照实验}
  \label{tab:robustness}
  \begin{tabularx}{\textwidth}{LCCC}
    \toprule
    变更项 & 变更内容 & 结果 & 是否翻转 \\
    \midrule
    换口径 & ... & ... & ... \\
    换种子 & ... & ... & ... \\
    换实现 & ... & ... & ... \\
    \bottomrule
  \end{tabularx}
\end{table}

\subsection{适用边界与失效域}

说明模型在何种参数区间、何种数据规模、何种假设下成立，以及越出边界后首先失效的是哪一条。凡在检验中被降级的结论，此处须写明其实际强度等级，并与降级声明逐条一致。


% ── 全文通用形态（Mrite 高教社杯规范）──
% 正文一律自然段；禁分点符号、禁 	extbf（摘要“针对问题X”与问题重述“问题N：”除外）。
% 跨页表用 longtable，短表用 tabularx；列宽比例总和 = 1.04 − 0.04 × N（N 为列数），格内文字不换行。
% 图宽 0.8	extwidth，表宽 	extwidth；图题交给 \caption，绘图代码不写 set_title()；图表必须在正文被引用。
% 章与章、节与节之间一律不加 
ewpage。
% ── 本节合同 ──
% 优点写“做到了什么 + 证据在哪”，不写空泛褒词。
% 缺点按“缺陷 — 影响 — 改进”三段写，且必须与降级声明一致：凡被降级的结论，此处要写明实际强度等级。
%
\section{模型的评价}

\subsection{模型的优点}
\begin{itemize}[itemindent=2em]
\item \textbf{优点1：}
\item \textbf{优点2：}
\item \textbf{优点3：}
\end{itemize}

\subsection{模型的缺点}
\begin{itemize}[itemindent=2em]
\item \textbf{缺点1：}
\item \textbf{缺点2：}
\end{itemize}


% ── 全文通用形态 ──
% 第8章只写模型改进；模型推广按2026参照模板删除。
% 本节的简要 AI 声明已经位于参考文献之前，不得另设重复声明。
% 详细记录只放支撑材料的 AI工具使用详情.pdf；正文中对 AI 生成内容的使用位置应作标注。

\section{模型的改进}

模型的改进应针对模型的缺点提出可操作方案，说明替换的方法、增加的数据或放宽的假设，并说明预期收益。

\subsection{模型的改进方向}

……（针对模型缺点逐条说明改进方向、实现方式和预期收益。）

\subsection*{AI工具使用声明}

% 在 main.tex 中仅当实际使用 AI 时设置 \aiusedtrue；不得填写模板示例为事实。
\ifaiused
本参赛队在竞赛过程中使用了AI工具，主要用于【简要用途，如语言润色、代码调试等】，详细使用情况见支撑材料中的《AI工具使用详情.pdf》。
\else
本参赛队在竞赛过程中未使用 AI 工具。
\fi



%%%%%%%%%%%%%%%%%%%%%%%%%%%%%%%%%%%%%%%%%%%%%%%%%%%%%%%%%%%%%
%% 参考文献
\input{9.参考文献.tex}


% ── 附录 ──
% \label{sec:appendix} 不得删除：latex_gate.py 依赖它切分正文页数与附录页数。
% 附录承载支撑材料索引、核心代码、环境和复现说明。

\makeatletter
\let\@oldaddcontentsline\addcontentsline
\renewcommand{\addcontentsline}[3]{}
\section*{附录}
\label{sec:appendix}
\let\addcontentsline\@oldaddcontentsline
\makeatother

\subsection*{附录1：支撑材料目录}
\label{app:support-files}

逐项列出随论文提交的支撑材料，路径与 delivery-manifest.json 保持一致。竞赛已提供的原始数据不重复提交；自行获得的数据才可在支撑材料中列出。

\begin{table}[H]
  \centering
  \caption{支撑材料目录}
  \label{tab:appendix_files}
  \renewcommand{\arraystretch}{1.2}
  \begin{tabularx}{\textwidth}{LXX}
    \toprule
    相对路径 & 类型 & 用途与论文位置 \\
    \midrule
    支撑材料/source/ & 完整源码 & 完整、可运行的项目；入口为支撑材料/source/main.py。 \\
    支撑材料/data/ & 自行获得数据 & 仅按实际需要提交，并写明来源与处理方式。 \\
    支撑材料/results/ & 补充结果 & 放正文容纳不下的结果、表格或图表。 \\
    \bottomrule
  \end{tabularx}
\end{table}

\subsection*{附录2：核心源程序节选}
\label{app:core-code}

本节只放说明模型逻辑所需的核心代码，不打印整个项目。完整可运行源码及依赖文件位于支撑材料/source/，实际入口为支撑材料/source/main.py。

\begin{verbatim}
# Replace with the model's core logic.
def solve(instance):
    # model construction, key constraint, and decision logic
    ...
\end{verbatim}

节选应对应正文的模型、约束或求解逻辑；不得把这一小段节选说成完整源码。

\subsection*{附录3：运行环境与复现说明}
\label{app:reproduction}

请如实写明语言、版本、依赖、操作系统条件、输入输出位置和复现命令。最终包根目录的 reproduce.py 应执行至少一条真实的最小重算命令，而不是只做文件存在性检查。

\begin{verbatim}
# Example only. Replace with the actual environment and command.
python --version
python reproduce.py
\end{verbatim}



\end{document}
